(1) 哈密顿量迹为零，特征方程 \(\det(H-EI)=E^2-[(vt)^2+\Delta^2]/4=0\)，瞬时本征能为 \(E_\pm(t)=\pm\frac12\sqrt{(vt)^2+\Delta^2}\)。两能级间隙 \(E_+-E_-=\sqrt{(vt)^2+\Delta^2}\)，在 \(t=0\) 取严格正的最小值 \(\Delta\)。
(2) 定义混合角 \(\tan2\vartheta=\Delta/(vt)\)，其分支按连续本征矢选择。远离交叉区 \(|vt|\gg\Delta\)，耦合相对对角能差极小，瞬时本征态近似哈密顿量的两个原始基态；穿过 \(t=0\) 时两条绝热能级避免真正交叉。
(3) 标准Landau–Zener精确散射结果给从下绝热能级跃入上绝热能级的概率 \(P_{\rm nonad}=\exp[-\pi\Delta^2/(2\hbar v)]\)。该概率对应时间从负无穷延伸至正无穷、参数线性变化且耦合常数固定的模型；不是单纯采用局部绝热定理即可无证明得到的恒等式。
(4) 由于 \([\Delta^2/(\hbar v)]=1\)，指数无量纲。\(v\to0^+\) 时指数趋负无穷，跃迁概率趋零而演化绝热；\(v\to\infty\) 时趋一；\(\Delta\to0^+\) 时避免交叉消失，固定原始态随时间延续会变成另一条绝热能级，故非绝热概率趋一。
